\section{Conclusion}\label{sec:conclusion}

% ------------------------------------------------------------------------------
\subsection{Synthesis of Key Findings}\label{sec:conclusion:synthesis}
% ------------------------------------------------------------------------------

Tracing the trajectory from Dennis Gabor's optical wave interference~\cite{GABOR_1948_HOLOGRAPHY} and John Sutherland's complex-valued Holographic Neural Networks~\cite{SUTHERLAND_1990_HNN}, through Tony Plate's Holographic Reduced Representations~\cite{PLATE_1995_HRR} and Pentti Kanerva's Hyperdimensional Computing~\cite{KANERVA_2009_HDC}, to contemporary Complex-Valued Neural Networks~\cite{BASSEY_CVNN_SURVEY}, Fourier token mixers~\cite{LEE_THORP_FNET}, linear attention~\cite{KATHAROPOULOS_2020_LINEAR_TRANSFORMER}, and State-Space Models~\cite{GU_SSM_S4_MAMBA} reveals a remarkably unified mathematical, architectural, and physical picture.
Three core takeaways synthesize the findings of this survey.

\paragraph{Mathematical Synthesis: Distributed Phase Superposition and the $\SNR \propto D/N$ Law.}
Across nearly eight decades of seemingly disparate formulations, holographic neural computation rests on a single algebraic core formalized in Section~\ref{sec:math}: mapping structured data into high-dimensional complex phasor spaces ($\C^D$, and ideally the unitary torus $\mathbb{T}^D$) where quasi-orthogonality supports exponentially many distinct symbols (Proposition~\ref{prop:quasi_ortho}), binding via circular convolution $\cconv$ or $\bigo{D}$ spectral Hadamard phase addition $\hadamard$ preserves dimensionality (Proposition~\ref{prop:conv_thm}), bundling superposes $N$ associations into a fixed-size trace, and circular correlation $\ccorr$ unbinds queried entries as the exact Hilbert-space adjoint (Lemma~\ref{lem:ccorr_props}).
This algebraic compression trades the $\bigo{ND}$ storage of explicit tables for $\bigo{1}$ memory with respect to $N$, incurring zero-mean additive crosstalk noise $\vect{\eta}$ governed universally by $\SNR = \Theta(D/N)$ (Theorem~\ref{thm:snr_capacity}).
Crucially, unitary complex phasors eliminate spectral self-distortion and reject orthogonal quadrature noise, securing a factor-of-two ($3\,\mathrm{dB}$) discrimination $\SNR$ advantage over real Gaussian HRR at equal width $D$, while clean-up projection onto a discrete codebook $\mathcal{V}$ of size $M$ yields reliable retrieval capacity $N_{\max} = \Theta(D / \ln(M/\delta))$ (Corollary~\ref{cor:cleanup_capacity}).
Rather than an arbitrary heuristic, the Shannon-entropy ceiling of fixed-width traces (Proposition~\ref{prop:superposition_duality}) provides an exact analytical boundary separating compressed semantic recall from lossless verbatim storage.

\paragraph{Architectural Synthesis: Complementarity Over Replacement.}
Our genealogical analysis (Section~\ref{sec:history}) and six-axis cross-paradigm evaluation (Section~\ref{sec:comparison}, Table~\ref{tab:paradigm_comparison}) establish a sharp structural bifurcation across modern AI.
\emph{Non-parametric explicit attention}---encompassing Modern Hopfield Networks~\cite{RAMSAUER_MODERN_HOPFIELD} and softmax Transformers~\cite{VASWANI_TRANSFORMER}---eliminates additive crosstalk via unfactorable exponential selectivity, achieving lossless verbatim copying and fine-grained disambiguation at the cost of $\bigo{ND}$ KV-cache growth and $\bigo{N^2}$ pairwise interaction.
\emph{Parametric holographic superposition}---spanning HNNs, HRR/HDC, linear attention, and selective SSMs---delivers constant-memory $\bigo{1}$ state updates, native dimension-preserving role--filler compositionality, one-shot algebraic plasticity, and graceful noise degradation, yet cannot circumvent fixed-state entropy bounds on unstructured verbatim copying~\cite{JELASSI_2024_REPEAT_AFTER_ME,ARORA_2023_ZOOLOGY}.
The decisive architectural conclusion is therefore \emph{complementarity rather than wholesale replacement}: holographic superposition is not a drop-in substitute for exact local attention, but an ideal mathematical substrate for compressed long-horizon semantic memory, structured neuro-symbolic binding, and interference-controlled continual adaptation.
As articulated in Section~\ref{sec:potentials}, hybrid foundation architectures (Figure~\ref{fig:hybrid_arch}) that pair bounded exact sliding-window attention with multi-timescale holographic $\bigo{1}$ memory modules, leverage crosstalk coherence diagnostics ($\rho_{\mathrm{coh}}$) for calibrated epistemic abstention, and multiplex task adapters across orthogonal complex phase subspaces offer a principled pathway toward scalable, verifiable, and lifelong AI.

\paragraph{Hardware Co-Design Synthesis: Matching Algebraic Primitives to Physical Substrates.}
Perhaps the clearest lesson of Section~\ref{sec:applications} (Table~\ref{tab:hardware_substrates}) is that holographic representations achieve their most transformative gains when paired with physical substrates aligned to their algebraic primitives rather than forced onto real-valued GEMM pipelines.
Discrete binary and bipolar VSA specializations collapse binding and similarity to bitwise XOR and popcount, enabling single-pass on-device learning on milliwatt digital microcontrollers and FPGAs~\cite{RAHIMI_HDC_EDGE,KIM_2020_GENIEHD}, while their high-dimensional bit-error resilience ($10\text{--}30\%$ BER tolerance) naturally absorbs conductance drift and stuck-at faults in analog Phase-Change Memory and RRAM in-memory computing crossbars~\cite{KARUNARATNE_2020_PCM_HDC}.
At the continuous end of the spectrum, complex-phasor HNNs and FHRR exhibit a one-to-one physical isomorphism with coherent optical wavefields: passive interference performs bundling, spatial phase masks execute Hadamard binding, spherical lenses compute exact spatial Fourier transforms, and integrated Mach--Zehnder interferometer meshes~\cite{SHEN_PHOTONIC_MZI,CLEMENTS_2016_MZI} alongside free-space diffractive networks ($\text{D}^2\text{NN}$)~\cite{LIN_DIFFRACTIVE_OPTICAL_NN} realize complex unitary transformations at optical carrier speeds~\cite{WETZSTEIN_2020_OPTICAL_AI,MCMAHON_2023_PHYSICS_OPTICAL_NN}.

% ------------------------------------------------------------------------------
\subsection{Outlook for Algorithm--Hardware Co-Design}\label{sec:conclusion:outlook}
% ------------------------------------------------------------------------------

Looking ahead, the convergence of hybrid holographic--Transformer/SSM algorithms with emerging photonic interferometers, diffractive optics, and analog in-memory fabrics signals a rare inflection point in computing.
For decades, algorithmic progress has been shaped by the CMOS GEMM \emph{hardware lottery}~\cite{HOOKER_2021_HARDWARE_LOTTERY}, privileging dense real-valued matrix multiplications that execute efficiently on digital systolic arrays while penalizing complex phase operations, FFT convolutions, and sparse algebraic bindings.
As electronic scaling encounters hard physical ceilings in interconnect energy and off-chip memory bandwidth, computational primitives that operate natively in complex phase space, tolerate analog physical noise, and maintain constant-footprint distributed state are poised to move from niche academic alternatives to first-class architectural building blocks.

Realizing this promise at frontier scale will require coordinated progress across three concrete research priorities:
\begin{enumerate}
    \item \textbf{Scalable Complex-Valued Training Primitives and Hybrid Memory Architectures:}
    The deep learning systems community must overcome phase-wrapping optimization instabilities in deep phase-coupled networks through principled manifold optimization, phase-preserving normalization, and stable Wirtinger recipes~\cite{BASSEY_CVNN_SURVEY,TRABELSI_2018_DCN}. In parallel, bridging the immediate digital hardware gap demands fused GPU/accelerator kernels (analogous to FlashAttention and parallel selective scans) for batched complex FFTs, gated Hadamard recurrences, and hierarchical sliding-window plus holographic $\bigo{1}$ KV-cache blocks (Section~\ref{sec:potentials:kvcache}).
    \item \textbf{Rigorous Capacity, Correlation, and Uncertainty Theory for Learned Non-Orthogonal Keys:}
    Classical VSA capacity theorems assume i.i.d.\ random quasi-orthogonal keys, whereas end-to-end neural encoders inevitably induce semantic clustering and cross-key correlations that bias superposed readouts. Extending the analytical framework of Section~\ref{sec:math} to correlated learned keys---combining mutual-coherence regularization, block-sparse phase codes, resonator/Hopfield clean-up attractors~\cite{FRADY_2020_RESONATOR,RAMSAUER_MODERN_HOPFIELD}, and calibrated $\SNR$-based epistemic uncertainty bounds ($\rho_{\mathrm{coh}}$)---is essential to unlock both high-capacity long-context memory and verifiable, hallucination-resistant neuro-symbolic reasoning (Sections~\ref{sec:potentials:reasoning}--\ref{sec:potentials:continual}).
    \item \textbf{Co-Designed Optical, Analog, and Edge Benchmarks Bridging Theory and Physical Deployment:}
    On the hardware frontier, closing the gap between laboratory photonic prototypes and general-purpose AI accelerators requires amortizing electro-optic conversion (E--O--E) overheads via high-fan-in spatial and wavelength multiplexing, maturing low-power optical nonlinearities, and leveraging non-volatile photonic phase-change materials to eliminate static thermal drift~\cite{WETZSTEIN_2020_OPTICAL_AI,MCMAHON_2023_PHYSICS_OPTICAL_NN}. Simultaneously, the field needs standardized, compute- and energy-normalized cross-paradigm benchmarks spanning long-horizon associative recall, compositional systematicity, continual task retention, and physical noise resilience (Section~\ref{sec:potentials:challenges}).
\end{enumerate}

The journey from optical interference plates to high-dimensional latent superposition has come full circle: mathematical abstractions originally borrowed from wave physics to explain biological associative memory now stand ready to both alleviate the scaling bottlenecks of modern foundation models and guide the architecture of future optical and analog AI hardware.
By uniting rigorous phase-space theory, hybrid neural--symbolic architectures, and physically native execution substrates, holographic neural computation offers a durable blueprint for more energy-efficient, compositional, and resilient artificial intelligence.
